\documentclass[12pt,a4paper]{article}

\usepackage[utf8]{inputenc}
\usepackage[T2A]{fontenc}
\usepackage[russian]{babel}

\usepackage{amsmath,amssymb,amsthm}
\usepackage{geometry}
\geometry{margin=2.2cm}
\usepackage{microtype}
\usepackage{parskip}

\newcommand{\Z}{\mathbb{Z}}
\DeclareMathOperator{\gcdop}{gcd}
\DeclareMathOperator{\ord}{ord}
\DeclareMathOperator{\Ker}{Ker}

\begin{document}

\section*{\centering Лекция 3}

\subsection*{Цикличность $(\Z/p)^{\ast}$, первообразные корни}

\textbf{Факт.} Для простого $p$ группа $(\Z/p)^{\ast}$ циклическая.

\textbf{Определение.} Элемент $g\in(\Z/p)^{\ast}$ называется \textbf{первообразным корнем по модулю $p$}, если
\[
\ord_p(g)=p-1,
\]
то есть $\langle g\rangle=(\Z/p)^{\ast}$.

\textbf{Число первообразных корней.}
Если $(\Z/p)^{\ast}=\langle g\rangle$, то каждый элемент имеет вид $g^m$, $m\in\{0,1,\dots,p-2\}$.
Элемент $g^m$ является образующим тогда и только тогда, когда
\[
\gcdop(m,p-1)=1.
\]
Следовательно, количество первообразных корней равно $\varphi(p-1)$.

\textbf{Изоморфизм с аддитивной группой.}
Отображение
\[
\psi:(\Z/(p-1),+)\to(\Z/p)^{\ast},\qquad \psi(\overline{m})=g^m
\]
— изоморфизм групп (при фиксированном первообразном корне $g$).

\medskip
\textbf{Классификация цикличности $(\Z/n\Z)^{\ast}$.}
Группа $(\Z/n\Z)^{\ast}$ циклическая тогда и только тогда, когда
\[
n\in\{2,\;4,\;p^s,\;2p^s\},
\]
где $p$ — нечётное простое, $s\in\mathbb{N}$.

\subsection*{Лемма Гензеля (поднятие корня)}

\textbf{Лемма Гензеля.}
Пусть $f(x)\in\Z[x]$, $p$ — простое, $j\ge 1$.
Пусть $a\in\Z$ таково, что
\[
f(a)\equiv 0\pmod{p^j}
\quad\text{и}\quad
f'(a)\not\equiv 0\pmod p.
\]
Тогда существует единственное $\tilde a$ по модулю $p^{j+1}$ такое, что
\[
\tilde a\equiv a\pmod{p^j},
\qquad
f(\tilde a)\equiv 0\pmod{p^{j+1}}.
\]
Это $\tilde a$ называется \textbf{поднятием} решения $a$.

\textbf{Доказательство (идея через разложение Тейлора).}
Ищем $\tilde a$ в виде
\[
\tilde a=a+t p^j,\qquad t\in\Z.
\]
По формуле Тейлора (с остатком по модулю $p^{j+1}$) имеем
\[
f(a+t p^j)\equiv f(a)+t p^j f'(a)\pmod{p^{j+1}},
\]
потому что все члены с $(t p^j)^2$ делятся на $p^{2j}$, а при $j\ge 1$ это кратно $p^{j+1}$.

Нужно добиться $f(a+t p^j)\equiv 0\pmod{p^{j+1}}$, то есть
\[
f(a)+t p^j f'(a)\equiv 0\pmod{p^{j+1}}.
\]
Так как $f(a)\equiv 0\pmod{p^j}$, можно разделить на $p^j$:
\[
\frac{f(a)}{p^j}+t f'(a)\equiv 0\pmod p,
\]
или
\[
t f'(a)\equiv -\frac{f(a)}{p^j}\pmod p.
\]
По условию $f'(a)\not\equiv 0\pmod p$, значит $f'(a)$ обратим по модулю $p$, и решение $t$ по модулю $p$ существует и единственно.
Следовательно, существует и единственное $\tilde a=a+t p^j \ (\mathrm{mod}\ p^{j+1})$.

\subsection*{Первообразные корни по модулю $p^2$ и $p^r$ при $p\ge 3$}

Далее $p$ — нечётное простое.

\subsection*{Переход от $p$ к $p^2$}

\textbf{Утверждение.} Если $g$ — первообразный корень по модулю $p$, то существует $t\in\{0,1,\dots,p-1\}$ такой, что
\[
g+t p
\]
— первообразный корень по модулю $p^2$.

\textbf{Доказательство.}
Рассмотрим все элементы вида $g+t p$ по модулю $p^2$.
Так как $\gcdop(g,p)=1$, то и $\gcdop(g+t p,p)=1$, значит $g+t p\in(\Z/p^2\Z)^{\ast}$.

Обозначим $h_t=\ord_{p^2}(g+t p)$. Тогда
\[
h_t \mid \varphi(p^2)=p(p-1).
\]
Кроме того, по модулю $p$ имеем $g+t p \equiv g$, значит
\[
(g+t p)^{p-1}\equiv g^{p-1}\equiv 1\pmod p.
\]
Следовательно, порядок $h_t$ кратен $p-1$:
\[
p-1 \mid h_t.
\]
Итак, возможны только два случая:
\[
h_t=p-1 \quad\text{или}\quad h_t=p(p-1).
\]

Рассмотрим многочлен
\[
f(x)=x^{p-1}-1.
\]
Поскольку $g$ — первообразный корень по модулю $p$, имеем $g^{p-1}\equiv 1\pmod p$, то есть $f(g)\equiv 0\pmod p$.
Производная
\[
f'(x)=(p-1)x^{p-2},
\]
и
\[
f'(g)\equiv (p-1)g^{p-2}\not\equiv 0\pmod p,
\]
так как $p\nmid(p-1)$ и $g$ обратим по модулю $p$.

По лемме Гензеля существует единственное поднятие $\tilde g$ по модулю $p^2$ такое, что
\[
\tilde g\equiv g\pmod p,
\qquad
\tilde g^{p-1}\equiv 1\pmod{p^2}.
\]
Иными словами, \emph{существует ровно один} представитель вида $g+t p$, для которого
\[
(g+t p)^{p-1}\equiv 1\pmod{p^2}.
\]
Для этого единственного $t$ порядок $h_t$ делит $p-1$, значит $h_t=p-1$.

Для всех остальных $t$ выполняется
\[
(g+t p)^{p-1}\not\equiv 1\pmod{p^2},
\]
а поскольку $h_t$ кратен $p-1$ и делит $p(p-1)$, остаётся единственная возможность:
\[
h_t=p(p-1).
\]
Значит, для некоторого $t$ элемент $g+t p$ имеет порядок $p(p-1)$ и является первообразным корнем по модулю $p^2$.

\subsection*{Индукция по степени: первообразный корень по модулю $p^r$}

\textbf{Утверждение.} Для каждого $r\ge 1$ группа $(\Z/p^r\Z)^{\ast}$ циклическая, и существует элемент $g$ такой, что
\[
\ord_{p^r}(g)=\varphi(p^r)=p^{r-1}(p-1).
\]

\textbf{Идея индукции.}
Пусть для некоторого $r\ge 2$ уже найден $g$ такой, что
\[
\ord_{p^r}(g)=p^{r-1}(p-1).
\]
Тогда, в частности,
\[
g^{p^{r-2}(p-1)}\equiv 1\pmod{p^{r-1}}
\quad\text{и}\quad
g^{p^{r-2}(p-1)}\not\equiv 1\pmod{p^{r}}.
\]
Запишем
\[
g^{p^{r-2}(p-1)} \equiv 1+t p^{r-1}\pmod{p^{r}},
\qquad t\not\equiv 0\pmod p.
\]
Возведём обе части в степень $p$:
\[
g^{p^{r-1}(p-1)} \equiv (1+t p^{r-1})^{p}\pmod{p^{r+1}}.
\]
По биному Ньютона:
\[
(1+t p^{r-1})^{p}
=1+\binom{p}{1}t p^{r-1}+\binom{p}{2}t^2 p^{2(r-1)}+\cdots
\equiv 1+p\cdot t p^{r-1}\pmod{p^{r+1}},
\]
поскольку все члены начиная с $\binom{p}{2}$ содержат множитель $p^2$ и степень $p^{2(r-1)}$, что при $r\ge 2$ даёт кратность $p^{r+1}$.
Следовательно,
\[
g^{p^{r-1}(p-1)} \equiv 1+t p^{r}\not\equiv 1\pmod{p^{r+1}}.
\]
Это и есть ключевой шаг, показывающий, что порядок на следующей степени действительно умножается на $p$.

\subsection*{Следствия: число первообразных корней и случай $2p^r$}

\textbf{Следствие 1.} Число первообразных корней по модулю $p^r$ равно
\[
\varphi(\varphi(p^r))=\varphi\bigl(p^{r-1}(p-1)\bigr)=p^{r-2}(p-1)\,\varphi(p-1)\quad (r\ge 2).
\]
(Для $r=1$ получаем $\varphi(p-1)$.)

\textbf{Следствие 2.} Для $p$ нечётного и $r\ge 1$ имеем изоморфизм
\[
(\Z/2p^r)^{\ast}\cong(\Z/p^r)^{\ast},
\]
поэтому и число первообразных корней по модулю $2p^r$ совпадает с числом первообразных корней по модулю $p^r$.
*
\section{Тест Миллера--Рабина}

Пусть $n>2$ — нечётное число. Представим
\[
n-1=2^{s}q,\qquad q\ \text{нечётно}.
\]

\subsection*{Наблюдение для простого модуля}

\textbf{Лемма-наблюдение.}
Если $p$ — простое и $p-1=2^s q$ ($q$ нечётно), то для любого $a\in(\Z/p)^{\ast}$ выполнено одно из двух:
\begin{enumerate}
  \item $a^q\equiv 1\pmod p$;
  \item существует $i\in\{0,1,\dots,s-1\}$ такое, что
  \[
  a^{2^{i}q}\equiv -1\pmod p.
  \]
\end{enumerate}

\textbf{Доказательство.}
Рассмотрим последовательность
\[
a^{q},\ a^{2q},\ a^{4q},\ \dots,\ a^{2^{s}q}=a^{p-1}\equiv 1\pmod p.
\]
Пусть $i$ — минимальный индекс такой, что $a^{2^{i}q}\equiv 1\pmod p$.
Если $i=0$, то $a^{q}\equiv 1\pmod p$ и выполнен пункт (1).
Если $i>0$, то положим $x=a^{2^{i-1}q}$. Тогда
\[
x^2=a^{2^{i}q}\equiv 1\pmod p,
\]
то есть $(x-1)(x+1)\equiv 0\pmod p$.
Так как $p$ простое, отсюда $x\equiv 1\pmod p$ или $x\equiv -1\pmod p$.
Но $i$ минимален, значит $x\not\equiv 1\pmod p$, поэтому $x\equiv -1\pmod p$, что и даёт пункт (2).

\subsection*{Формулировка теста}

\textbf{Тест Миллера--Рабина для нечётного $n$.}
Берём некоторое основание $a$ со взаимной простотой $\gcdop(a,n)=1$ и проверяем:
\begin{enumerate}
  \item вычисляем $a^{q}\pmod n$;
  \item если $a^{q}\equiv 1\pmod n$, то $n$ проходит проверку для основания $a$;
  \item иначе последовательно возводим в квадрат:
  \[
  a^{q},\ a^{2q},\ a^{4q},\dots,a^{2^{s-1}q}\pmod n
  \]
  и если на каком-то шаге получаем $-1\pmod n$, то $n$ проходит проверку для основания $a$ - свидетель простоты;
  \item иначе $a$ является \textbf{свидетелем непростоты} и $n$ составно.
\end{enumerate}

\subsection{*Оценка Рабина: не менее $\tfrac34$ оснований обнаруживают составность}

Пусть $m$ — нечётное составное число, и
\[
m-1=2^{\ell}q,\qquad q\ \text{нечётно}.
\]
Рассмотрим множество оснований, для которых $m$ проходит тест Миллера--Рабина (т.\,е. «обманывает» тест).

\textbf{Теорема Рабина, идея}
Если $m$ составно, то доля оснований $a\in(\Z/m)^{\ast}$, на которых $m$ проходит тест Миллера--Рабина, не превосходит $\frac14$.

\subsection*{Сведение через КТО и гомоморфизм}

Пусть разложение на простые множители имеет вид
\[
m=p_1p_2\cdots p_n,\qquad n\ge 3,
\]
и рассмотрим изоморфизм по китайской теореме об остатках:
\[
\Phi:(\Z/m)^{\ast}\xrightarrow{\ \sim\ }(\Z/p_1\Z)^{\ast}\times\cdots\times(\Z/p_n\Z)^{\ast},
\quad
\Phi(a)=(a\bmod p_1,\dots,a\bmod p_n).
\]

Для каждого $i$ разложим
\[
p_i-1=2^{k_i}q_i,\qquad q_i\ \text{нечётно},
\]
и положим
\[
\ell'=\min\{k_1,\dots,k_n\}.
\]
Рассмотрим гомоморфизм
\[
\alpha:(\Z/m)^{\ast}\to(\Z/m)^{\ast},\qquad \alpha(a)=a^{2^{\ell'}q}.
\]
Тогда $\Ker(\alpha)$ — подгруппа, и все прообразы элементов из $\Im(\alpha)$ имеют одинаковую мощность, равную $|\Ker(\alpha)|$.

Определим множества:
\[
G=\{a\in(\Z/m)^{\ast}\mid a^{2^{\ell'}q}\equiv 1\pmod m\},
\]
\[
H=\{a\in(\Z/m)^{\ast}\mid a^{2^{\ell'}q}\equiv \pm 1\pmod m\}.
\]
Ясно, что множество «лжецов» теста Миллера--Рабина содержится в $H$.

\subsection*{Подсчёт мощностей $G$ и $H$}

В произведении групп по КТО условие $a^{2^{\ell'}q}\equiv 1\pmod m$ означает, что каждая координата равна $1$:
\[
\Phi\bigl(\alpha(a)\bigr)=(1,1,\dots,1).
\]
Следовательно,
\[
G=\alpha^{-1}\bigl(\Phi^{-1}(1,1,\dots,1)\bigr).
\]
Это означает, что $G$ является одним слоем (одним прообразом) гомоморфизма $\alpha$:
\[
|G|=|\Ker(\alpha)|.
\]
Однако аналогично можно рассматривать и все кортежи вида $(\pm 1,\pm 1,\dots,\pm 1)$.
Всего таких кортежей $2^{n}$, и каждому соответствует прообраз при $\alpha$, имеющий мощность $|\Ker(\alpha)|$.
Тогда общее число элементов, у которых $\alpha(a)$ в точности равна некоторому кортежу из $\{\pm 1\}^n$, равно
\[
2^{n}\cdot |\Ker(\alpha)|.
\]

Из этих кортежей ровно два являются «однородными»:
\[
(1,1,\dots,1)\quad\text{и}\quad(-1,-1,\dots,-1).
\]
Им соответствуют элементы из $H$, поэтому
\[
|H|=2\cdot |\Ker(\alpha)|.
\]
А множеству всех кортежей из $\{\pm 1\}^n$ соответствует множество размера
\[
2^{n}\cdot |\Ker(\alpha)|.
\]
Отсюда отношение:
\[
\frac{|H|}{2^{n}\cdot |\Ker(\alpha)|}=\frac{2}{2^{n}}=\frac{1}{2^{n-1}}\le \frac14\quad (n\ge 3).
\]
Значит, доля «лжецов» ограничена сверху $\frac14$.

\textbf{Вывод о вероятности ошибки.}
Если выполнить тест $t$ раз на независимых основаниях, то вероятность того, что составное $m$ пройдет все проверки, не превосходит
\[
\left(\frac14\right)^{t}.
\]

\textbf{Замечание (эвристики).}
Существуют эвристические утверждения о том, что для проверки простоты достаточно перебрать основания $a$ до величин порядка $O(\log^2 m)$, но это отдельный вопрос.

\textbf{Гипотеза Артина}
Существует бесконечно много простых $p$, для которых $2$ является первообразным корнем по модулю $p$.

\subsection*{Упражнение на дом про RSA}

Пусть
\[
N=pq,
\]
где $p$ и $q$ — большие простые.

Берём элемент $g\in(\Z/N\Z)^{\ast}$ и выбираем параметры $r_1,r_2$.
Определим
\[
g_1 \equiv g^{r_1(p-1)}\pmod N,
\qquad
g_2 \equiv g^{r_2(q-1)}\pmod N.
\]
Для сообщения $m\in\Z/N\Z$ формируем пару значений, например:
\[
c_1 \equiv m\cdot g_1^{s_1}\pmod N,
\qquad
c_2 \equiv m\cdot g_2^{s_2}\pmod N,
\]
где $s_1,s_2$ — дополнительные параметры (в зависимости от конкретной схемы).

\textbf{Ключевое свойство (сведение по модулю простых).}
По малой теореме Ферма для любого $x\in(\Z/p)^{\ast}$ выполнено $x^{p-1}\equiv 1\pmod p$.
Следовательно,
\[
g_1 \equiv g^{r_1(p-1)}\equiv 1\pmod p,
\qquad
g_2 \equiv g^{r_2(q-1)}\equiv 1\pmod q.
\]
Отсюда
\[
c_1 \equiv m\pmod p,
\qquad
c_2 \equiv m\pmod q,
\]
(при подходящем выборе/нейтрализации степеней $s_1,s_2$).
Зная $p$ и $q$, можно восстановить $m$ из системы
\[
\begin{cases}
x\equiv c_1\pmod p,\\
x\equiv c_2\pmod q,
\end{cases}
\]
по китайской теореме об остатках, получая единственное решение $x\equiv m\pmod N$.

\textbf{Замечание о факторизации через $\gcd$.}
Так как $g_1\equiv 1\pmod p$, то $p\mid (g_1-1)$, и потому
\[
\gcdop(N,g_1-1)=p.
\]
Аналогично,
\[
\gcdop(N,g_2-1)=q.
\]
Это объясняет, почему элементы с «встроенной» конгруэнцией $1$ по одному из простых модулей напрямую связаны с задачей факторизации.

\subsection*{Квадратичные вычеты}

Пусть $p$ — нечётное простое.

\textbf{Определение.}
Элемент $a\in(\Z/p)^{\ast}$ называется \textbf{квадратичным вычетом} по модулю $p$, если существует $x$ такое, что
\[
x^2\equiv a\pmod p.
\]

\textbf{Критерий Эйлера.}
Для $a\in(\Z/p)^{\ast}$ верно:
\[
a\ \text{— квадрат по модулю }p
\iff
a^{\frac{p-1}{2}}\equiv 1\pmod p.
\]

\textbf{Доказательство.}
Так как $(\Z/p)^{\ast}$ циклическая, пусть $(\Z/p)^{\ast}=\langle g\rangle$.
Тогда $a=g^{m}$ для некоторого $m$.

Элемент $a$ является квадратом тогда и только тогда, когда существует $y$ с
\[
y^2\equiv a\pmod p.
\]
Пусть $y=g^{t}$. Тогда
\[
y^2=g^{2t}\equiv g^m \pmod p \iff 2t\equiv m\pmod{p-1}.
\]
Последнее сравнение разрешимо тогда и только тогда, когда $m$ чётно.

С другой стороны,
\[
a^{\frac{p-1}{2}}=(g^m)^{\frac{p-1}{2}}=g^{m\frac{p-1}{2}}.
\]
Так как $g^{\frac{p-1}{2}}\equiv -1\pmod p$ (элемент порядка $2$), получаем
\[
g^{m\frac{p-1}{2}}\equiv 1\pmod p \iff m\ \text{чётно}.
\]
Итак, $a$ — квадрат $\iff a^{\frac{p-1}{2}}\equiv 1\pmod p$.

\subsubsection*{Гомоморфизм и подсчёт числа квадратов}

Рассмотрим отображение
\[
\alpha:(\Z/p)^{\ast}\to(\Z/p)^{\ast},\qquad \alpha(a)=a^{\frac{p-1}{2}}.
\]
Это гомоморфизм.

\textbf{Ядро.}
\[
\Ker(\alpha)=\left\{a\in(\Z/p)^{\ast}\mid a^{\frac{p-1}{2}}\equiv 1\pmod p\right\}
\]
— в точности множество квадратичных вычетов.

\textbf{Образ.}
Так как $a^{p-1}\equiv 1\pmod p$, имеем
\[
\alpha(a)^2=a^{p-1}\equiv 1\pmod p,
\]
значит $\alpha(a)\in\{1,-1\}$ и
\[
\Im(\alpha)\subseteq\{1,-1\}.
\]
На деле $\Im(\alpha)=\{1,-1\}$, потому что $\alpha(1)=1$, а для первообразного корня $g$ получаем
\[
\alpha(g)=g^{\frac{p-1}{2}}\equiv -1\pmod p.
\]

По теореме о гомоморфизме
\[
|(\Z/p)^{\ast}|=|\Ker(\alpha)|\cdot |\Im(\alpha)|.
\]
Левая часть равна $p-1$, а $|\Im(\alpha)|=2$, значит
\[
|\Ker(\alpha)|=\frac{p-1}{2}.
\]
То есть квадратов (квадратичных вычетов) ровно $\frac{p-1}{2}$.

\subsection*{Пример теоремы Гензеля (не получилось)}
\subsubsection*{попытка применить Гензеля для сравнения по модулю $27$}

Рассмотрим сравнение
\[
x^4+x^2+6x+1\equiv 0\pmod{27}.
\]
Обозначим
\[
f(x)=x^4+x^2+6x+1,\qquad f'(x)=4x^3+2x+6.
\]

\subsubsection*{Шаг 1: модуль $3$}
По модулю $3$ имеем $6x\equiv 0$, поэтому
\[
f(x)\equiv x^4+x^2+1\pmod 3.
\]
Проверка даёт решения
\[
x\equiv 1\pmod 3,\qquad x\equiv 2\pmod 3.
\]
Проверим условие Гензеля $f'(a)\not\equiv 0\pmod 3$:
\[
f'(1)=4+2+6=12\equiv 0\pmod 3,
\]
\[
f'(2)=4\cdot 8+4+6=42\equiv 0\pmod 3.
\]
Условие не выполняется ни для одного корня по модулю $3$, поэтому прямое применение леммы Гензеля с $p=3$ невозможно.

\subsubsection*{Шаг 2: модуль $9$}
Если искать решения по модулю $9$, получаются кандидаты
\[
x\in\{1,4,7\}\pmod 9.
\]
Однако проверка производной по модулю $3$ снова даёт нуль:
\[
f'(1)=12\equiv 0\pmod 3,\qquad
f'(4)=4\cdot 64+8+6=270\equiv 0\pmod 3,\qquad
f'(7)=4\cdot 343+14+6=1392\equiv 0\pmod 3.
\]
Следовательно, даже при работе по модулю $9$ условие $f'(a)\not\equiv 0\pmod 3$ не выполняется, и поднять решения до модуля $27$ методом Гензеля в этой форме не удаётся.

\subsection*{Немного про изоморфизм}

\textbf{Пример.} Так как $24=2^3\cdot 3$, то
\[
\Z/24 \cong \Z/8 \times \Z/3.
\]
Но при этом
\[
\Z/24 \not\cong \Z/6 \times \Z/4,
\]
потому что $6$ и $4$ не взаимно просты, порядок группы слева не достигает 24.

\textbf{Пример.} Диэдральная группа $D_3$ (симметрии правильного треугольника) изоморфна $S_3$.
Это видно из того, что действие симметрий треугольника на множестве его вершин задаёт вложение $D_3\hookrightarrow S_3$, и обе группы имеют порядок $6$.

\textbf{Замечание.} Группа $S_4$ не коммутативна. Ввиду этого рассматривать изоморфизм этой группы для $\Z/24$ не имееет смысла.

\end{document}